
\begin{corollaire}
\label{XIII.5.6}
Sean $X$ un $S$-esquema y
$D=\sum_{1\leqslant i\leqslant r}\divisor f_i$ un divisor de
cruzamientos normales relativamente a~$S$, como en~\Ref{XIII.5.4}.
Sean $\overline{x}$ un punto geométrico de~$X$, $\overline{X}$ la
localización estricta de~$X$ en~$\overline{x}$,
$\overline{Y}=Y_{(\overline{X})}$,
$\overline{U}=\overline{X}-\overline{Y}$ y
$$
\widetilde{U}=\varprojlim_{(n_i)}
\overline{U}[T_i]_{i\in I(x)}\big/(T_i^{n_i}-f_i)\quoi,
$$
donde el límite proyectivo se toma sobre el conjunto dirigido de las
familias de enteros $n_i>0$, primos con la característica~$p$
de~$\kres(x)$. Entonces $\widetilde{U}$ es un recubrimiento universal
moderadamente ramificado de~$\overline{U}$ relativamente a~$S$. Por
consiguiente, el grupo fundamental moderadamente ramificado
de~$\overline{U}$ es
$$
\pi_1^{\tame}(\overline{U})\simeq
\prod_{\ell\neq p}\ZZ_\ell[1]^{I(x)}
\qquad \textrm{(isomorfismo canónico)\quoi.}
$$
El grupo $\pi_1^{\tame}(\overline{U})$ es isomorfo, de manera no
canónica, a~$\prod_{\ell\neq p}\ZZ_\ell^{I(x)}$.
\end{corollaire}

\setcounter{subsection}{6}
\setcounter{subsubsection}{0}
\begin{subremarque}
\label{XIII.5.6.1}
Sean $X$ un $S$-esquema y
$D=\sum_{1\leqslant i\leqslant r}\divisor f_i$ un divisor de
cruzamientos normales relativamente a~$S$, como en~\Ref{XIII.5.4}, y
$U=X-\Supp D$. Para todo subconjunto $I\subset[1,r]$, sea
$$
X_I=\Big(\tbigcap_{i\in I}V(f_i)\Big)\cap
\Big(\tbigcap_{i\in\complement I}X_{f_i}\Big)\quoi.
$$
Sea $p$ un número primo o cero, y sea $Z$ un subconjunto de~$X_I$
todos cuyos puntos son de característica~$p$. Sea
$$
\widetilde{U}_I=\varprojlim_{(n_i)}
\ifthenelse{\boolean{orig}}
{U[T_i]}
{U[T_i]_{i\in I}}
\big/(T_i^{n_i}-f_i)\quoi,
$$
donde
\marginpar{434}
el límite proyectivo se toma sobre el conjunto dirigido de las familias
de enteros $n_i>0$, primos con~$p$. Entonces, para todo punto geométrico
$\overline{x}$ de~$Z$, la imagen inversa de~$\widetilde{U}_I$ sobre
$\overline{U}$ se identifica con el recubrimiento universal moderadamente
ramificado de~$\overline{U}$.
\end{subremarque}
